% !TEX root = ctfp-print.tex

\lettrine[lhang=0.17]{T}{his is the Scala edition} of \emph{Category Theory for Programmers}.
It's been a tremendous success, making Bartosz Milewski's blog post series available as a nicely 
typeset \acronym{PDF}, as well as a hardcover book. There have been numerous contributions made
to improve the book, by fixing typos and errors, as well as translating the code snippets into
other programming languages.

I am thrilled to present this edition of the book, containing the original Haskell code, followed by
its Scala counterpart. The Scala code snippets were generously provided by 
\urlref{https://github.com/typelevel/CT_from_Programmers.scala}{Typelevel} contributors, slightly 
modified to suit the format of this book.

To support code snippets in multiple languages, I am using a \LaTeX{} macro to load the code snippets
from external files. This allows easily extending the book with other languages, while leaving the 
original text intact. Which is why you should mentally append the words ``and Scala'' whenver you see
``in Haskell'' in the text.

The code is laid out in the following manner: the original Haskell snippet, followed by a Scala snippet.
To distinguish between them, the snippets are braced from the left with a vertical bar, in the primary
color of the language's logo, \raisebox{-.2mm}{\includegraphics[height=.3cm]{fig/icons/haskell.png}}, 
and \raisebox{-.2mm}{\includegraphics[height=.3cm]{fig/icons/scala.png}} respectively, e.g.:

\srcsnippet{content/3.6/code/haskell/snippet03.hs}{blue}{haskell}
\unskip
\srcsnippet{content/3.6/code/scala/snippet03.scala}{red}{scala}
\NoIndentAfterThis

In addition, some Scala snippets make use of the 
\urlref{https://github.com/non/kind-projector}{Kind Projector} compiler plugin, to support nicer syntax for partially-applied types.